\subsection{Биномиальные коэффициенты}
\textbf{Базовое.} $\binom nk=\binom n{n-k}=\frac nk\binom{n-1}{k-1}$; Паскаль $\binom nk=\binom{n-1}{k-1}+\binom{n-1}k$; $\binom{-n}{k}=(-1)^k\binom{n+k-1}k$.
\par\textbf{Суммы по $k$.} $\sum_k\binom nk=2^n$; чётные и нечётные $k$ дают по $2^{n-1}$; $\sum_k k\binom nk=n2^{n-1}$;
$\sum_k\binom nk\binom km=\binom nm2^{n-m}$; $\sum_k\binom nk^2=\binom{2n}n$; $\sum_{k}\binom{n-k}k=F_{n+1}$.
\par $\sum_{i\le k}(-1)^i\binom ni=(-1)^k\binom{n-1}k$.
\par\textbf{Вандермонд:} $\sum_k\binom mk\binom n{r-k}=\binom{m+n}r$.
\par\textbf{Клюшка:} $\sum_{i=r}^n\binom ir=\binom{n+1}{r+1}$;\ \ диагональ $\sum_{i\le k}\binom{n+i}i=\binom{n+k+1}k$.
\par $\sum_{i=0}^k\binom ki\big/\binom ni=\binom{n+1}{n-k+1}\big/\binom nk$.
\par $\sum_i\binom ni i^k=\sum_jS(k,j)\,n^{\underline j}\,2^{n-j}$.
\par\textbf{Скользящая сумма} $B(l,r,n)=\sum_{i=l}^r\binom ni$: $B(l,r,n{+}1)=2B(l,r,n)+\binom n{l-1}-\binom nr$ (сдвиг $n$ за $O(1)$).
\par\textbf{Куммер:} $v_p\binom nm$ = число переносов при сложении $m$ и $n-m$ в $p$-ичной системе. \textbf{Лукас:} $\binom nm\equiv\prod_i\binom{n_i}{m_i}\pmod p$ ($n_i,m_i$ -- $p$-ичные цифры).
\subsection{Разложение на слагаемые}
$x_1+\dots+x_n=m$: при $x_i\ge0$ -- $\binom{m+n-1}{n-1}$; при $x_i\ge1$ -- $\binom{m-1}{n-1}$;
при $0\le x_i\le k$ -- $\sum_{t=0}^{\min(n,\lfloor m/(k+1)\rfloor)}(-1)^t\binom nt\binom{m-t(k+1)+n-1}{n-1}$.
\par $n$ различимых людей по $k$ различимым группам размера $\le m$ (пустые можно): $f(0,k)=1$, $f(n,k)=k\,f(n{-}1,k)-k\binom{n-1}m f(n{-}m{-}1,k{-}1)$.
\par Пары путей $(0,0)\to(n,m)$ по сетке без общих точек, кроме концов (неупорядоченные): $\binom{n+m-2}{n-1}^2-\binom{n+m-2}n\binom{n+m-2}m$. Общий случай: Линдстрем--Гессель--Вьенно (det матрицы чисел путей). Пути, задевающие прямую: отражение начала.
\subsection{Каталан и скобочные последовательности}
$C_n=\frac1{n+1}\binom{2n}n=\binom{2n}n-\binom{2n}{n-1}=\frac1{2n+1}\binom{2n+1}n$,\ \ $C_n=\sum_{i<n}C_iC_{n-1-i}$.\\
{\scriptsize 1, 1, 2, 5, 14, 42, 132, 429, 1430, 4862, 16796, 58786, 208012, 742900, 2674440}.
Считает: ПСП, триангуляции $(n{+}2)$-угольника, бинарные деревья на $n$ вершинах, пути под диагональю, непересекающиеся хорды.
$\sum_{x_0+\dots+x_k=n}\prod C_{x_i}=\frac{k+1}{n+k+1}\binom{2n+k}n$.
\par \textbf{ПСП:} баланс всех префиксов $\ge0$, итог $0$. «Дополняемых до ПСП» с $n$ открывающими и $m\le n$ закрывающими: $\binom{n+m}m-\binom{n+m}{m-1}$. Ровно $k$ пиков «()» (числа Нараяны): $N(n,k)=\frac1n\binom nk\binom n{k-1}$; столько же плоских корневых деревьев с $n$ рёбрами и $k$ листьями. Скобка с номером $p$ (с нуля) закрывающая в $\sum_{m=0}^{\min(n-1,\lfloor(p-1)/2\rfloor)}C_mC_{n-m-1}$ ПСП. Идея вывода: $2n{+}1$ циклических сдвигов слова из $n$ «(» и $n{+}1$ «)» ровно один начинается с ПСП.
\subsection{Числа Стирлинга, Белла, Эйлера}
$S(n,k)$ -- разбиения $n$-множества на $k$ блоков: $S(n,k)=kS(n{-}1,k)+S(n{-}1,k{-}1)=\frac1{k!}\sum_j(-1)^{k-j}\binom kjj^n$; $x^n=\sum_kS(n,k)x^{\underline k}$; $S(n,k)=[x^{n-k}]\prod_{i=1}^k\frac1{1-ix}$; вся строка: $S(n,k)=[x^k]\big(\sum_i\frac{i^n}{i!}x^i\big)\big(\sum_j\frac{(-1)^j}{j!}x^j\big)$ (одна свёртка).
\par $c(n,k)$ -- перестановки с $k$ циклами (беззнаковые первого рода): $c(n,k)=c(n{-}1,k{-}1)+(n{-}1)c(n{-}1,k)$; $x^{\overline n}=\sum_kc(n,k)x^k$, $x^{\underline n}=\sum_k(-1)^{n-k}c(n,k)x^k$; столбец: $\frac{(-\ln(1-x))^k}{k!}$, строка: $\prod_{t<n}(x+t)$ (D\&C + FFT).
\par \textbf{Обращение Стирлинга:} $P(k)=\sum_iS(k,i)Q(i)\iff Q(k)=\sum_i(-1)^{k-i}c(k,i)P(i)$.
\par Блоки с разностью соседних элементов $\ge d$: $S(n{-}d{+}1,\,k{-}d{+}1)$ при $k\ge d$.\ \ Белл: $B_n=\sum_kS(n,k)=\sum_{k}\binom{n-1}{k-1}B_{n-k}$; {\scriptsize 1, 1, 2, 5, 15, 52, 203, 877, 4140, 21147, 115975, 678570, 4213597}.
\par \textbf{Эйлеровы} $E(n,k)$ (перестановок с $k$ подъёмами): $E(n,k)=(n{-}k)E(n{-}1,k{-}1)+(k{+}1)E(n{-}1,k)=\sum_j(-1)^j\binom{n+1}j(k{+}1{-}j)^n$; $\sum_kE(n,k)=n!$.
\subsection{Перестановки}
\textbf{Беспорядки:} $D_n=(n{-}1)(D_{n-1}+D_{n-2})=nD_{n-1}+(-1)^n=n!\sum_{t\le n}\frac{(-1)^t}{t!}=\lfloor n!/e\rceil$ ($n\ge1$), $\sum_k\binom nkD_{n-k}=n!$; {\scriptsize 1, 0, 1, 2, 9, 44, 265, 1854, 14833, 133496}.
\par Среднее по случайной перестановке: циклов -- $H_n$, рекордов (префиксных максимумов) -- $H_n$, неподвижных точек -- $1$.
\par Циклы длин из множества $T$: $\sum_n g_T(n)\frac{x^n}{n!}=\exp\big(\sum_{n\in T}\frac{x^n}n\big)$.
\par \textbf{Бернсайд:} классов $=\frac1{|G|}\sum_{g}|X^g|$. Ожерелья из $n$ бусин в $k$ цветах с точностью до поворота: $\frac1n\sum_{d\mid n}\varphi(d)k^{n/d}$.
\par Бинарных строк ровно с $m$ различными подпоследовательностями (считая пустую): $\varphi(m{+}1)$.
\subsection{Разбиения числа}
$\prod_{n\ge1}(1-x^n)=\sum_{k\in\mathbb Z}(-1)^kx^{k(3k-1)/2}$ (пентагональная теорема), отсюда $p(n)=\sum_{k\ne0}(-1)^{k+1}p\big(n-\tfrac{k(3k-1)}2\big)$ за $O(n\sqrt n)$. $p(n)\approx\frac1{4\sqrt3n}e^{\pi\sqrt{2n/3}}\approx\frac{0.145}ne^{2.56\sqrt n}$.\\
{\scriptsize $p(0..10)=1,1,2,3,5,7,11,15,22,30,42$;\ $p(20)=627$,\ $p(50)=204226$,\ $p(100)=190569292$}.
\subsection{Деревья, остовы, эйлеровы циклы}
\textbf{Кэли:} $n^{n-2}$ помеченных деревьев; код Прюфера -- биекция с последовательностями из $n{-}2$ чисел. С заданными степенями $d_v$ ($\sum d_v=2n{-}2$): $\frac{(n-2)!}{\prod(d_v-1)!}$. Лесов на $n$ вершинах из $m$ деревьев, где вершины $1..m$ в разных деревьях: $m\,n^{n-m-1}$\,$^\dagger$. $\tau(K_{a,b})=a^{b-1}b^{a-1}$.
\par \textbf{Кирхгоф:} для ребра $u\to v$: $L[u][v]{-}{=}1,\ L[v][v]{+}{=}1$ (неориентированное -- оба направления). Минор без строки и столбца $r$ -- число остовных деревьев, направленных \emph{от} $r$ (любой $r$ для неориентированного). С $L[u][u]{+}{=}1$ вместо $L[v][v]$ -- деревья, направленные \emph{к} корню.
\par \textbf{BEST:} число эйлеровых циклов орграфа (как циклических последовательностей рёбер) $=t_w(G)\prod_v(\deg_{out}(v){-}1)!$, где $t_w$ -- число остовов, направленных к $w$\,$^\dagger$.
\par \textbf{Эйлеров путь:} в неориентированном -- связность по рёбрам и $\le2$ нечётных вершин; в орграфе -- $\mathrm{out}-\mathrm{in}\in\{0,\pm1\}$, не более одной «начальной» и одной «конечной».
\subsection{Производящие функции (помеченные объекты, EGF)}
$\mathrm{Seq}(A)=\frac1{1-A}$, $\mathrm{Set}(A)=e^{A}$, $\mathrm{Cycle}(A)=\ln\frac1{1-A}$, пара $=A\cdot B$. Следствия: $S(n,k)\leftrightarrow\frac{(e^x-1)^k}{k!}$, $c(n,k)\leftrightarrow\frac{(-\ln(1-x))^k}{k!}$, $D_n\leftrightarrow\frac{e^{-x}}{1-x}$, Белл $\leftrightarrow e^{e^x-1}$. Непомеченные: $\mathrm{MSet}(A)=\exp\big(\sum_{k\ge1}\frac{A(z^k)}k\big)$, $\mathrm{Cycle}(A)=\sum_{n\ge1}\frac{\varphi(n)}n\ln\frac1{1-A(z^n)}$.
\subsection{Суммы и ряды}
$\sum_{i\le n} i^2=\frac{n(n{+}1)(2n{+}1)}6$,\ $\sum i^3=\frac{n^2(n{+}1)^2}4$,\ $\sum i^4=\frac{n(n{+}1)(2n{+}1)(3n^2{+}3n{-}1)}{30}$;\ \ $\sum_{i=1}^nia^i=\frac{a(na^{n+1}-(n{+}1)a^n+1)}{(a-1)^2}$.
\par $e^x=\sum\frac{x^k}{k!}$, $\ln(1{+}x)=\sum_{k\ge1}\frac{(-1)^{k+1}x^k}k$, $(1{-}x)^{-n}=\sum_k\binom{n+k-1}{n-1}x^k$, $\sin,\cos$ -- знакочередующиеся нечётные/чётные; $\arctan x=\sum(-1)^k\frac{x^{2k+1}}{2k+1}$. Если $a_n=\sum_{i\le k}c_ia_{n-i}$ и корни $r_j$ характеристического многочлена различны, то $a_n=\sum d_jr_j^n$ (кратный корень даёт множитель-многочлен).
